\documentclass[aspectratio=169, 11pt]{beamer}

\usetheme{Madrid}
\usecolortheme{seahorse}
\setbeamertemplate{navigation symbols}{}
\setbeamertemplate{footline}[frame number]

\usepackage[T1]{fontenc}
\usepackage[utf8]{inputenc}
\usepackage[french]{babel}
\usepackage{amsmath, amssymb}
\usepackage{booktabs}
\usepackage{tikz}
\usetikzlibrary{shapes, arrows.meta, positioning, calc}
\usepackage{multicol}

% ── Couleurs ─────────────────────────────────────────────────────────────────
\definecolor{accent}{HTML}{2196F3}
\definecolor{fair}{HTML}{4CAF50}
\definecolor{alertc}{HTML}{d9534f}
\definecolor{neutral}{HTML}{FF9800}

% ── Titre ────────────────────────────────────────────────────────────────────
\title{Parcours p\'edagogique : IA Responsable}
\subtitle{Du jeu de donn\'ees aux conclusions --- pas \`a pas}
\author{IA708 --- T\'el\'ecom Paris}
\date{2026}

\begin{document}

% ══════════════════════════════════════════════════════════════════════════════
\begin{frame}
\titlepage
\end{frame}

\begin{frame}{Plan}
\tableofcontents
\end{frame}

% ══════════════════════════════════════════════════════════════════════════════
\section{Vue d'ensemble}

\begin{frame}{Le parcours en 11 \'etapes}
\centering
\begin{tikzpicture}[
    node distance=0.35cm and 0.8cm,
    block/.style={rectangle, draw, rounded corners=3pt, minimum height=0.75cm,
                  minimum width=2cm, text centered, font=\scriptsize,
                  text width=1.9cm, align=center},
    question/.style={font=\tiny\itshape, text=accent!80!black, text width=2.5cm, align=center},
    arrow/.style={-Stealth, thick, accent!70}
]
\node[block, fill=accent!12] (data) {1.~Donn\'ees\newline 1000 lignes};
\node[block, fill=accent!12, right=of data] (explore) {2.~Exploration\newline biais?};
\node[block, fill=accent!12, right=of explore] (prep) {3.~Pr\'eparation\newline split, encoding};
\node[block, fill=accent!12, right=of prep] (model) {4.~Mod\`ele\newline r\'eg.\ logist.};

\node[block, fill=fair!12, below=0.9cm of data] (metrics) {5.~M\'etriques\newline perf + \'equit\'e};
\node[block, fill=fair!12, right=of metrics] (rw) {6.~Reweighing\newline pr\'e-trait.};
\node[block, fill=fair!12, right=of rw] (pp) {7.~Seuils/grp\newline post-trait.};
\node[block, fill=fair!12, right=of pp] (comp) {8.~Comparer\newline 4 configs};

\node[block, fill=neutral!12, below=0.9cm of metrics] (shap) {9.~SHAP\newline interpr\'et.};
\node[block, fill=neutral!12, right=of shap] (rob) {10.~Robustesse\newline perturbation};
\node[block, fill=alertc!12, right=of rob] (ccl) {11.~Conclusion\newline le\c{c}ons};

\draw[arrow] (data) -- (explore);
\draw[arrow] (explore) -- (prep);
\draw[arrow] (prep) -- (model);
\draw[arrow] (model.south) -- ++(0,-0.25) -| (metrics.north);
\draw[arrow] (metrics) -- (rw);
\draw[arrow] (rw) -- (pp);
\draw[arrow] (pp) -- (comp);
\draw[arrow] (comp.south) -- ++(0,-0.25) -| (shap.north);
\draw[arrow] (shap) -- (rob);
\draw[arrow] (rob) -- (ccl);

\node[question, above=0.1cm of explore] {Y a-t-il des biais ?};
\node[question, above=0.1cm of model] {Quelles perf.\ ?};
\node[question, below=0.1cm of rw] {Peut-on corriger ?};
\node[question, below=0.1cm of comp] {Quel compromis ?};
\node[question, below=0.1cm of shap] {Pourquoi ces d\'ecisions ?};
\node[question, below=0.1cm of ccl] {Qu'a-t-on appris ?};
\end{tikzpicture}

\vspace{0.6em}
\small Chaque \'etape r\'epond \`a une question qui motive la suivante.
Tout cod\'e \emph{from scratch} (numpy/pandas).
\end{frame}

% ══════════════════════════════════════════════════════════════════════════════
\section{Les donn\'ees}

\begin{frame}{\'Etape 1 --- Le German Credit Dataset}
\begin{columns}[T]
\begin{column}{0.55\textwidth}
\textbf{UCI Statlog} (Hofmann, 1994)
\begin{itemize}
    \item 1\,000 demandes de cr\'edit
    \item 20 attributs (7 num., 13 cat.)
    \item Cible : d\'efaut de paiement (0/1)
\end{itemize}

\vspace{0.5em}
\textbf{Pourquoi ce dataset ?}
\begin{itemize}
    \item Scoring cr\'edit = d\'ecision \`a fort impact
    \item Risque de discrimination historique
    \item Cas d'\'ecole en IA responsable
\end{itemize}
\end{column}
\begin{column}{0.42\textwidth}
\centering\small
\begin{tabular}{lc}
\toprule
\textbf{Stat} & \textbf{Valeur} \\
\midrule
Bon (0) & 700 (70\%) \\
Mauvais (1) & 300 (30\%) \\
\midrule
Num\'eriques & 7 \\
Cat\'egorielles & 13 \\
\bottomrule
\end{tabular}

\vspace{0.8em}
\begin{alertblock}{}
\small D\'es\'equilibre 70/30 $\Rightarrow$\\ accuracy na\"ive trompeuse
\end{alertblock}
\end{column}
\end{columns}
\end{frame}

% ══════════════════════════════════════════════════════════════════════════════
\section{Exploration}

\begin{frame}{\'Etape 2 --- Biais pr\'eexistants dans les donn\'ees}
\textbf{Question} : avant tout mod\`ele, y a-t-il d\'ej\`a des disparit\'es entre groupes ?

\vspace{0.5em}
\begin{columns}[T]
\begin{column}{0.48\textwidth}
\centering\textbf{\textcolor{accent}{Genre}}

\small
\begin{tabular}{lcc}
\toprule
& \textbf{Bon} & \textbf{D\'efaut} \\
\midrule
Homme (690) & 72\% & 28\% \\
Femme (310) & 65\% & 35\% \\
\bottomrule
\end{tabular}

\vspace{0.3em}
\scriptsize \'Ecart $\approx$ 7 points
\end{column}
\begin{column}{0.48\textwidth}
\centering\textbf{\textcolor{fair}{\^Age}}

\small
\begin{tabular}{lcc}
\toprule
& \textbf{Bon} & \textbf{D\'efaut} \\
\midrule
Adulte $>$25 (810) & 72\% & 28\% \\
Jeune $\leq$25 (190) & 60\% & 40\% \\
\bottomrule
\end{tabular}

\vspace{0.3em}
\scriptsize \'Ecart $\approx$ 12 points
\end{column}
\end{columns}

\vspace{0.8em}
\begin{alertblock}{Cons\'equence}
Un mod\`ele performant \textbf{reproduira} ces disparit\'es historiques.
Bonne performance $\neq$ mod\`ele \'equitable.
\end{alertblock}
\end{frame}

\begin{frame}{\'Etape 2 --- Deux attributs sensibles}
\begin{columns}[T]
\begin{column}{0.48\textwidth}
\textbf{\textcolor{accent}{Genre}}
\begin{itemize}
    \item Extrait de \texttt{personal\_status\_sex}
    \item A91/A93/A94 $\to$ homme
    \item A92/A95 $\to$ femme
    \item Privil\'egi\'e : homme
\end{itemize}
\end{column}
\begin{column}{0.48\textwidth}
\textbf{\textcolor{fair}{\^Age}}
\begin{itemize}
    \item Binaris\'e au seuil de 25 ans
    \item Adulte ($> 25$) / jeune ($\leq 25$)
    \item Privil\'egi\'e : adulte
    \item Biais plus marqu\'e (+12 pts)
\end{itemize}
\end{column}
\end{columns}

\vspace{0.8em}
\begin{block}{D\'ecision de conception}
Les deux colonnes sont \textbf{retir\'ees des features} $\Rightarrow$ pas de discrimination \emph{directe}.

\vspace{0.2em}
Mais des \textbf{proxies} (features corr\'el\'ees) peuvent encoder indirectement l'information sensible. \`A v\'erifier via SHAP (\'etape~9).
\end{block}
\end{frame}

% ══════════════════════════════════════════════════════════════════════════════
\section{Pr\'eparation}

\begin{frame}{\'Etape 3 --- Pr\'eparation des donn\'ees}
\begin{columns}[T]
\begin{column}{0.52\textwidth}
\textbf{Pr\'etraitement}
\begin{itemize}
    \item \textbf{Num\'eriques} (6 col.) : z-score
    $$x' = \frac{x - \mu_{\text{train}}}{\sigma_{\text{train}}}$$
    \item \textbf{Cat\'egorielles} (12 col.) : one-hot
    \item R\'esultat : $\sim$44 features
    \item Fit sur train, appliqu\'e au val/test
\end{itemize}
\end{column}
\begin{column}{0.45\textwidth}
\textbf{Split stratifi\'e 60/20/20}

\vspace{0.3em}
\small
\begin{tabular}{lcc}
\toprule
& \textbf{Taille} & \textbf{R\^ole} \\
\midrule
Train & $\sim$600 & Entra\^inement \\
Val & $\sim$200 & Seuil, PP \\
Test & $\sim$200 & \'Eval.\ finale \\
\bottomrule
\end{tabular}

\vspace{0.5em}
\scriptsize Seed fix\'e (42) pour reproductibilit\'e.
Stratifi\'e : chaque split conserve le ratio 70/30.
\end{column}
\end{columns}

\vspace{0.5em}
\begin{alertblock}{Attention --- pas de fuite d'information}
$\mu$ et $\sigma$ calcul\'es \textbf{uniquement sur le train} (ddof=1).
Les cat\'egories inconnues au val/test sont mises \`a z\'ero.
\end{alertblock}
\end{frame}

% ══════════════════════════════════════════════════════════════════════════════
\section{Mod\`ele}

\begin{frame}{\'Etape 4 --- R\'egression logistique from scratch}
\textbf{Mod\`ele} :
$P(\text{d\'efaut}=1 \mid x) = \sigma(w^\top x + b)$
\quad avec \quad
$\sigma(z) = \frac{1}{1 + e^{-z}}$

\vspace{0.5em}
\begin{columns}[T]
\begin{column}{0.48\textwidth}
\textbf{Entra\^inement}
\begin{itemize}
    \item Perte : cross-entropy + L2
    \item Adam ($\text{lr}=0.03$, $\beta_1\!=\!0.9$, $\beta_2\!=\!0.999$)
    \item R\'egularisation $\lambda = 0.01$
    \item Early stopping (patience 300)
\end{itemize}
\end{column}
\begin{column}{0.48\textwidth}
\textbf{Seuil de d\'ecision}
\begin{itemize}
    \item Pas fix\'e \`a 0.5 (d\'es\'equilibre !)
    \item Balayage $\tau \in [0.1,\; 0.9]$
    \item Maximise la \textbf{balanced accuracy} :
\end{itemize}
$$\text{BalAcc} = \tfrac{1}{2}(\text{TPR} + \text{TNR})$$
\end{column}
\end{columns}

\vspace{0.5em}
\begin{block}{Perte (avec poids pour le reweighing)}
$\mathcal{L} = -\frac{1}{n}\sum_{i} s_i \big[ y_i \ln p_i + (1-y_i) \ln(1 - p_i) \big] + \frac{\lambda}{2} \|w\|^2$
\hfill {\small ($s_i = 1$ pour le baseline)}
\end{block}
\end{frame}

% ══════════════════════════════════════════════════════════════════════════════
\section{M\'etriques}

\begin{frame}{\'Etape 5 --- M\'etriques de performance}
\textbf{Question} : comment \'evaluer un mod\`ele sur un dataset d\'es\'equilibr\'e ?

\vspace{0.5em}
\small
\begin{tabular}{lp{10cm}}
\toprule
\textbf{M\'etrique} & \textbf{Pourquoi ?} \\
\midrule
\textbf{AUC ROC} & Invariante au seuil. Capacit\'e \`a s\'eparer les classes. Calcul\'ee par Wilcoxon-Mann-Whitney. \\
\textbf{Balanced Acc.} & $\frac{1}{2}(\text{TPR}+\text{TNR})$. Adapt\'ee au d\'es\'equilibre (vs accuracy na\"ive). \\
\textbf{ECE} & Expected Calibration Error (cf.~cours incertitude). \\
\bottomrule
\end{tabular}

\vspace{0.5em}
\normalsize
\textbf{ECE} : $\displaystyle\text{ECE} = \sum_{b=1}^{B} \frac{|B_b|}{n}\, \big|\text{acc}(B_b) - \text{conf}(B_b)\big|$

\vspace{0.2em}
\small Mesure si la confiance du mod\`ele correspond \`a la r\'ealit\'e.
Un mod\`ele bien calibr\'e a ECE $\approx 0$.
\end{frame}

\begin{frame}{\'Etape 5 --- M\'etriques d'\'equit\'e (cf.~cours fairness)}
Soit $S$ l'attribut sensible, $Y$ la cible, $\hat{Y}$ la pr\'ediction.

\vspace{0.5em}
\begin{tabular}{lll}
\toprule
\textbf{Crit\`ere} & \textbf{D\'efinition} & \textbf{Signification} \\
\midrule
\textbf{Demographic Parity} & $\big|P(\hat{Y}{=}1|S{=}0) - P(\hat{Y}{=}1|S{=}1)\big|$ & M\^eme taux de s\'election \\
\textbf{Equal Opportunity} & $\big|\text{TPR}_{S=0} - \text{TPR}_{S=1}\big|$ & M\^eme TPR \\
\bottomrule
\end{tabular}

\vspace{0.8em}
\textbf{\'Evaluation du baseline} (test) :

\vspace{0.3em}
\small
\begin{columns}[T]
\begin{column}{0.48\textwidth}
\centering\textbf{Genre}
\begin{tabular}{lcc}
\toprule
& S\'election & TPR \\
\midrule
Homme & $\approx 0.41$ & $\approx 0.70$ \\
Femme & $\approx 0.50$ & $\approx 0.74$ \\
$|\Delta|$ & 0.094 & 0.044 \\
\bottomrule
\end{tabular}
\end{column}
\begin{column}{0.48\textwidth}
\centering\textbf{\^Age}
\begin{tabular}{lcc}
\toprule
& S\'election & TPR \\
\midrule
Adulte & $\approx 0.42$ & $\approx 0.77$ \\
Jeune & $\approx 0.51$ & $\approx 0.59$ \\
$|\Delta|$ & 0.081 & \textbf{0.179} \\
\bottomrule
\end{tabular}
\end{column}
\end{columns}

\vspace{0.5em}
\normalsize
$\Rightarrow$ Le baseline \textbf{n'est pas \'equitable}. |EO gap| pour l'\^age = 0.179 : \textbf{correction n\'ecessaire}.
\end{frame}

% ══════════════════════════════════════════════════════════════════════════════
\section{\'Equit\'e}

\begin{frame}{\'Etape 6 --- Reweighing (pr\'e-traitement)}
\textbf{Question} : peut-on corriger les biais en repondérant les exemples ?

\vspace{0.3em}
\textbf{Principe} (cf.~cours fairness-mitigation) : rendre $S \perp Y$ dans le train.

$$w_i = \frac{P(S{=}s_i) \cdot P(Y{=}y_i)}{P(S{=}s_i,\; Y{=}y_i)}$$

\begin{columns}[T]
\begin{column}{0.52\textwidth}
\textbf{Intuition} :
\begin{itemize}
    \item Sous-groupe sous-repr\'esent\'e dans une classe $\Rightarrow$ poids $w > 1$
    \item Surrepr\'esent\'e $\Rightarrow$ poids $w < 1$
    \item On entra\^ine un \textbf{nouveau mod\`ele} avec ces poids
    \item Un mod\`ele par attribut sensible
\end{itemize}
\end{column}
\begin{column}{0.44\textwidth}
\centering\small
\begin{tabular}{lcc}
\toprule
& $Y{=}0$ & $Y{=}1$ \\
\midrule
Privil\'egi\'e & $w \approx 1$ & $w > 1$ \\
Non-priv. & $w > 1$ & $w \approx 1$ \\
\bottomrule
\end{tabular}

\vspace{0.5em}
\scriptsize Direction typique des poids
\end{column}
\end{columns}
\end{frame}

\begin{frame}{\'Etape 7 --- Seuils par groupe (post-traitement)}
\textbf{Question} : peut-on corriger \textbf{sans r\'e-entra\^iner} ?

\vspace{0.5em}
\textbf{Principe} : ajuster le seuil \textbf{s\'epar\'ement par groupe} pour \'egaliser le taux de s\'election.

$$\tau_{\text{global}} \to \tau_{g_0}, \tau_{g_1} \quad\text{t.q.}\quad P(\hat{Y}{=}1|S{=}g_0) \approx P(\hat{Y}{=}1|S{=}g_1)$$

\begin{itemize}
    \item Seuils calibr\'es sur la \textbf{validation} ($\sim$200 exemples)
    \item Crit\`ere : Demographic Parity
\end{itemize}

\vspace{0.5em}
\begin{columns}[T]
\begin{column}{0.48\textwidth}
\begin{exampleblock}{Avantage}
Pas de r\'e-entra\^inement. M\^eme mod\`ele, seul le seuil change.
\end{exampleblock}
\end{column}
\begin{column}{0.48\textwidth}
\begin{alertblock}{Limite}
Peu de donn\'ees de calibration $\Rightarrow$ seuils potentiellement instables.
\end{alertblock}
\end{column}
\end{columns}
\end{frame}

\begin{frame}{\'Etape 8 --- Comparaison : 4 configurations}
\small
\begin{table}
\centering
\begin{tabular}{llcccc}
\toprule
\textbf{Attr.} & \textbf{Configuration} & \textbf{AUC} & \textbf{BalAcc} & $|\Delta\text{DP}|$ & $|\Delta\text{EO}|$ \\
\midrule
Genre & Baseline          & 0.793 & 0.701 & 0.094 & 0.044 \\
Genre & Reweighing        & 0.793 & 0.701 & 0.094 & 0.044 \\
Genre & Reweighing + PP   & 0.793 & 0.727 & 0.238 & 0.192 \\
\midrule
\^Age  & Baseline          & 0.793 & 0.701 & 0.081 & 0.179 \\
\^Age  & Reweighing        & 0.789 & 0.713 & 0.081 & 0.202 \\
\^Age  & Reweighing + PP   & 0.789 & 0.712 & \textbf{0.056} & 0.261 \\
\bottomrule
\end{tabular}
\end{table}

\vspace{0.3em}
\normalsize
\textbf{Observations cl\'es :}
\begin{itemize}
    \item \textbf{Genre} : reweighing \textbf{sans effet} (lignes identiques) $\Rightarrow$ pas de proxy
    \item \textbf{\^Age} : reweighing + PP r\'eduit $|\Delta\text{DP}|$ (0.081 $\to$ \textbf{0.056})
    \item \textbf{Mais} $|\Delta\text{EO}|$ augmente (0.179 $\to$ 0.261) $\Rightarrow$ \textbf{compromis}
\end{itemize}
\end{frame}

\begin{frame}{Le th\'eor\`eme d'impossibilit\'e en action}
\begin{alertblock}{Chouldechova 2017, Kleinberg 2016}
Il est \textbf{math\'ematiquement impossible} de satisfaire simultan\'ement Demographic Parity, Equal Opportunity et calibration, sauf si les taux de base sont \'egaux entre groupes.
\end{alertblock}

\vspace{0.5em}
\textbf{Ce qu'on observe} :

\begin{columns}[T]
\begin{column}{0.55\textwidth}
Pour l'\textbf{\^age} :
\begin{itemize}
    \item R\'eduire $|\Delta\text{DP}|$ $\Rightarrow$ $|\Delta\text{EO}|$ \textbf{augmente}
    \item 0.081 $\to$ 0.056 (DP \checkmark)
    \item 0.179 $\to$ 0.261 (EO $\times$)
\end{itemize}

\vspace{0.3em}
Ce n'est pas un \'echec de la m\'ethode --- c'est une \textbf{limite fondamentale}.
\end{column}
\begin{column}{0.42\textwidth}
\centering
\begin{tikzpicture}[
    node distance=1cm,
    crit/.style={rectangle, draw, rounded corners, font=\small\bfseries,
                 minimum width=1.8cm, minimum height=0.7cm, align=center}
]
\node[crit, fill=fair!15] (dp) {DP};
\node[crit, fill=accent!15, right=1cm of dp] (eo) {EO};
\coordinate (mid) at ($(dp)!0.5!(eo)$);
\node[crit, fill=neutral!15, below=0.8cm of mid] (cal) {Calibration};

\draw[<->, thick, alertc] (dp) -- node[above, font=\scriptsize] {tension} (eo);
\draw[<->, thick, alertc] (dp) -- (cal);
\draw[<->, thick, alertc] (eo) -- (cal);
\end{tikzpicture}

\vspace{0.3em}
\scriptsize On ne peut pas tout avoir.\\Il faut \textbf{choisir}.
\end{column}
\end{columns}
\end{frame}

% ══════════════════════════════════════════════════════════════════════════════
\section{Interpr\'etabilit\'e}

\begin{frame}{\'Etape 9 --- SHAP lin\'eaire exact}
\textbf{Question} : quelles features d\'eterminent les d\'ecisions ?

\vspace{0.3em}
Pour un mod\`ele lin\'eaire, les \textbf{valeurs de Shapley} sont exactes (pas d'approximation) :

$$\phi_i(x) = w_i \cdot \left(x_i - \mathbb{E}_{\text{train}}[x_i]\right)$$

\begin{columns}[T]
\begin{column}{0.48\textwidth}
\textbf{Top 5 features} (baseline) :
\begin{enumerate}
    \item \texttt{checking\_status}
    \item \texttt{duration\_in\_month}
    \item \texttt{credit\_amount}
    \item \texttt{credit\_history}
    \item \texttt{savings\_account\_bonds}
\end{enumerate}
\end{column}
\begin{column}{0.48\textwidth}
\textbf{Importance par permutation}
\begin{itemize}
    \item M\'elanger une colonne $\to$ mesurer la \textbf{chute d'AUC}
    \item 10 r\'ep\'etitions par feature
    \item M\'ethode compl\'ementaire au SHAP
\end{itemize}

\vspace{0.3em}
\begin{exampleblock}{}
\small Les deux m\'ethodes \textbf{convergent} sur le classement.
\end{exampleblock}
\end{column}
\end{columns}
\end{frame}

\begin{frame}{\'Etape 9 --- D\'etection de proxies}
\textbf{Question} : des features encodent-elles \textbf{indirectement} le genre ou l'\^age ?

\vspace{0.3em}
Un \textbf{proxy} = feature \`a la fois \textbf{corr\'el\'ee} ($|r| > 0.15$) \`a $S$ \textbf{et influente} (SHAP \'elev\'e).

\vspace{0.5em}
\begin{columns}[T]
\begin{column}{0.48\textwidth}
\textbf{\textcolor{accent}{Genre}}
\begin{itemize}
    \item Aucune feature \`a la fois influente ET corr\'el\'ee
    \item \textbf{Proxies faibles}
    \item $\Rightarrow$ Retirer la colonne suffit
    \item Coh\'erent : reweighing sans effet
\end{itemize}
\end{column}
\begin{column}{0.48\textwidth}
\textbf{\textcolor{fair}{\^Age}}
\begin{itemize}
    \item \texttt{duration\_in\_month} :\\ $|r| > 0.15$ + top 2 SHAP
    \item \textbf{Proxy actif}
    \item $\Rightarrow$ Le mod\`ele ``voit'' l'\^age
    \item Explique l'effet du reweighing
\end{itemize}
\end{column}
\end{columns}

\vspace{0.5em}
\begin{block}{Le\c{c}on fondamentale}
L'interpr\'etabilit\'e \textbf{explique} pourquoi le reweighing fonctionne pour l'\^age mais pas pour le genre. Sans SHAP, on ne comprendrait pas.
\end{block}
\end{frame}

% ══════════════════════════════════════════════════════════════════════════════
\section{Robustesse}

\begin{frame}{\'Etape 10 --- Perturbation et calibration}
\textbf{Question} : les corrections d'\'equit\'e fragilisent-elles le mod\`ele ?

\vspace{0.5em}
\begin{columns}[T]
\begin{column}{0.48\textwidth}
\textbf{Protocole}
\begin{itemize}
    \item Num. : $x_j + \varepsilon$, $\varepsilon \sim \mathcal{N}(0,\, 0.2 \cdot \sigma_j)$
    \item Cat. : swap al\'eatoire ($p{=}0.1$)
\end{itemize}

\vspace{0.3em}
\textbf{ECE} (cours incertitude) :
$$\text{ECE} = \sum_{b} \frac{|B_b|}{n} |\text{acc}_b - \text{conf}_b|$$

\small Mesure l'\'ecart confiance / r\'ealit\'e.
\end{column}
\begin{column}{0.48\textwidth}
\textbf{R\'esultats}

\small
\begin{tabular}{lccc}
\toprule
& AUC & BalAcc & ECE \\
\midrule
Clean & 0.793 & 0.701 & $\sim$0.05 \\
Perturb\'e & $\sim$0.77 & $\sim$0.68 & $\sim$0.06 \\
$\Delta$ & --0.02 & --0.02 & +0.01 \\
\bottomrule
\end{tabular}

\vspace{0.3em}
\normalsize
\begin{itemize}
    \item Corr\'elation scores : $r \approx 0.97$
    \item Mod\`eles fair \textbf{aussi robustes}
\end{itemize}
\end{column}
\end{columns}

\vspace{0.3em}
\begin{exampleblock}{Conclusion}
Les corrections d'\'equit\'e \textbf{n'amplifient pas la fragilit\'e}. Le mod\`ele lin\'eaire est naturellement robuste.
\end{exampleblock}
\end{frame}

% ══════════════════════════════════════════════════════════════════════════════
\section{Synth\`ese}

\begin{frame}{\'Etape 11 --- Ce qu'on a appris}
\begin{enumerate}
    \item \textbf{Les donn\'ees sont biais\'ees} --- taux de d\'efaut diff\'erents entre groupes.
          Un bon mod\`ele reproduira ces biais.

    \item \textbf{Retirer l'attribut sensible ne suffit pas toujours} ---
          genre : oui (proxies faibles).
          \^Age : non (\texttt{duration\_in\_month} = proxy actif).

    \item \textbf{Le reweighing d\'epend des proxies} ---
          sans proxy, aucun effet. Avec proxy, le mod\`ele change.

    \item \textbf{Compromis DP/EO in\'evitable} ---
          $|\Delta\text{DP}| \downarrow$ $\Rightarrow$ $|\Delta\text{EO}| \uparrow$
          (th\'eor\`eme d'impossibilit\'e).

    \item \textbf{Interpr\'etabilit\'e \'eclaire l'\'equit\'e} ---
          SHAP + corr\'elation = d\'etection de proxies.

    \item \textbf{Mod\`ele robuste} --- $r \approx 0.97$, m\^eme apr\`es corrections.
\end{enumerate}
\end{frame}

\begin{frame}{Les trois piliers sont li\'es}
\centering
\begin{tikzpicture}[
    node distance=1.8cm,
    pillar/.style={rectangle, draw, rounded corners=5pt, minimum width=3.5cm,
                   minimum height=1.3cm, font=\large\bfseries, text centered, align=center},
    link/.style={<->, very thick, gray!50}
]
\node[pillar, fill=fair!20] (eq) {\'Equit\'e};
\node[pillar, fill=neutral!20, right=3cm of eq] (interp) {Interpr\'etabilit\'e};
\node[pillar, fill=accent!20, below right=1.6cm and 1.5cm of eq] (rob) {Robustesse};

\draw[link] (eq) -- node[above, font=\small, text=gray!70!black] {proxies} (interp);
\draw[link] (eq) -- node[left, font=\small, text=gray!70!black, pos=0.4] {stabilit\'e} (rob);
\draw[link] (interp) -- node[right, font=\small, text=gray!70!black, pos=0.4] {confiance} (rob);
\end{tikzpicture}

\vspace{1em}
\begin{itemize}
    \item L'\textbf{interpr\'etabilit\'e} r\'ev\`ele les probl\`emes d'\textbf{\'equit\'e} (proxies)
    \item La \textbf{robustesse} garantit que les corrections d'\'equit\'e sont \textbf{fiables}
    \item La \textbf{confiance} dans le mod\`ele requiert les trois \`a la fois
\end{itemize}

\vspace{0.5em}
\centering\small\textit{Code from scratch --- numpy et pandas uniquement.}
\end{frame}

\end{document}
